\documentclass[11pt]{article}

\usepackage{amsmath, amssymb, bm}
\usepackage{geometry}
\geometry{margin=2.5cm}

\title{NeuroBridge: Shared Latent Dynamics, Spike Generation, and Representation Learning}
\author{Neuro\_Bridge}
\date{}

\begin{document}

\maketitle

\section{Purpose}

NeuroBridge is a controlled framework for studying whether latent task
structure can be recovered from neural population activity. The central idea is
to define a known low-dimensional task dynamics, generate subject-specific
spike-count observations from it, and then test whether representation-learning
methods recover both the latent geometry and the temporal relationship between
subjects.

\begin{equation}
\text{motor condition}
\rightarrow Z
\rightarrow U
\rightarrow \lambda
\rightarrow X
\end{equation}

where:

\begin{itemize}
    \item $Z$ is the low-dimensional latent trajectory.
    \item $U$ is the continuous neural drive.
    \item $\lambda$ is the positive firing-rate intensity.
    \item $X$ is the observed spike-count tensor.
\end{itemize}

\section{Notation and Dimensions}

Before specifying the model, we define the main mathematical objects.

\begin{center}
\begin{tabular}{lll}
\hline
Symbol & Meaning & General shape\\
\hline
$n_{\text{trials}}$ & number of repeated trials & scalar\\
$L$ & number of time bins in one trial & scalar\\
$\Delta t$ & duration of one bin & scalar\\
$k$ & latent dimension & scalar\\
$N$ & number of neurons & scalar\\
$c_i$ & condition/direction of trial $i$ & scalar\\
$Z_i(t)$ & shared task latent state & $k$-vector\\
$B^{(s)}$ & loading matrix for subject $s$ & $k \times N$\\
$u_i^{(s)}(t)$ & neural drive & $N$-vector\\
$\lambda_i^{(s)}(t)$ & firing-rate intensity & $N$-vector\\
$X_i^{(s)}(t)$ & spike-count vector & $N$-vector\\
$W_i^{(s)}(t)$ & neural window given to an encoder & window $\times N$\\
$E_i^{(s)}(t)$ & learned embedding & embedding vector\\
\hline
\end{tabular}
\end{center}

The full tensors are:

\begin{align}
Z_{\text{task}} &\in
\mathbb{R}^{n_{\text{trials}} \times L \times k},\\
X^{(s)} &\in
\mathbb{R}^{n_{\text{trials}} \times L \times N},\\
B^{(s)} &\in
\mathbb{R}^{k \times N},\\
E^{(s)} &\in
\mathbb{R}^{n_{\text{windows}} \times d_{\text{emb}}}.
\end{align}

Specific numerical values used in the poster run are reported later in the
experimental-configuration section.

\section{Latent Trajectories}

The latent tensor has shape:

\begin{equation}
Z_{\text{task}} \in
\mathbb{R}^{n_{\text{trials}} \times L \times k}.
\end{equation}

For each trial $i$:

\begin{equation}
Z_i(t) = m_i(t) + \eta_i(t).
\end{equation}

The deterministic part $m_i(t)$ defines the expected task trajectory for a
given condition. The stochastic part $\eta_i(t)$ introduces smooth
trial-to-trial variability around that expected trajectory.

\subsection{Deterministic Motor Component}

For a circular task, each condition is associated with an angle:

\begin{equation}
\theta_i = \frac{2\pi (c_i-1)}{n_{\text{conditions}}},
\end{equation}

where $c_i$ is the condition label for trial $i$.

The temporal movement profile is:

\begin{equation}
s(t) = 10t^3 - 15t^4 + 6t^5,
\end{equation}

with $t \in [0,1]$.

The polynomial $s(t)$ is a smooth phase variable for the trial. It maps the
normalized trial time from 0 to 1 while keeping the endpoints flat:
$s(0)=0$, $s(1)=1$, $s'(0)=s'(1)=0$, and
$s''(0)=s''(1)=0$. In this simulator it is not meant to be a full
biomechanical model of reaching; it is a controlled way to impose a smooth
within-trial progression.

The deterministic latent trajectory is:

\begin{equation}
m_i(t) =
\begin{bmatrix}
s(t)\cos\theta_i\\
s(t)\sin\theta_i\\
s(t)
\end{bmatrix}.
\end{equation}

Therefore, the first two coordinates describe directional displacement on a
circle, while the third coordinate stores the within-trial movement phase. This
third coordinate is important: it prevents the latent trajectory from being
only a static direction vector and gives the model an explicit temporal axis.

\subsection{Stochastic Component}

The noise term follows an AR(1)-like process:

\begin{equation}
\eta_i(t) = \phi \eta_i(t-1) + \epsilon_i(t),
\end{equation}

where:

\begin{equation}
\epsilon_i(t) \sim \mathcal{N}(0, \sigma_\epsilon^2 I).
\end{equation}

The innovation variance is $\sigma_\epsilon^2$. If the process is stationary
and $|\phi|<1$, the marginal variance of one latent-noise coordinate is:

\begin{equation}
\operatorname{Var}[\eta(t)] =
\frac{\sigma_\epsilon^2}{1-\phi^2}.
\end{equation}

Thus $\phi$ controls temporal persistence, while $\sigma_\epsilon$ controls the
amount of trial-to-trial fluctuation injected at each time step.

\section{Loading Matrix}

The loading matrix specifies how the latent state is read out by a neural
population. It is the bridge between task geometry and neuron-specific firing
modulation:

\begin{equation}
B \in \mathbb{R}^{k \times N}.
\end{equation}

Each column of $B$ corresponds to one neuron:

\begin{equation}
\bm{b}_j = B_{:,j} \in \mathbb{R}^{k}.
\end{equation}

\subsection{Directional Tuning}

The first two rows of $B$ define preferred directions:

\begin{align}
B_{1,j} &= a \cos \varphi_j,\\
B_{2,j} &= a \sin \varphi_j,
\end{align}

where $a$ is the directional scale.

Neurons are grouped by preferred task condition. This means the population is
structured: subsets of neurons preferentially respond to different directions.
The directional scale controls how strongly the circular latent coordinates
modulate the neural population. A larger value produces sharper directional
tuning and stronger separation between movement directions.

The remaining latent dimensions, if present, receive smaller random loadings:

\begin{equation}
B_{3:k,j} \sim \alpha_{\text{extra}} \mathcal{N}(0,1).
\end{equation}

The parameter $\alpha_{\text{extra}}$ controls weak sensitivity to
non-directional latent coordinates. In the circular setting, it lets the
movement-phase dimension influence firing rates without dominating the
directional geometry.

\subsection{Interpretation}

The matrix $B$ is not learned from data. It is imposed as part of the
generative model.

Its purpose is to define how neurons read out the latent motor geometry.
If a latent direction is aligned with the preferred direction of a neuron,
that neuron receives higher drive.

\section{Neural Drive}

The continuous neural drive is:

\begin{equation}
U^{(s)} = Z^{(s)} B^{(s)} + \bm{c}^{(s)}.
\end{equation}

With dimensions:

\begin{align}
Z^{(s)} &\in
\mathbb{R}^{n_{\text{trials}} \times L \times k},\\
B^{(s)} &\in \mathbb{R}^{k \times N},\\
\bm{c}^{(s)} &\in \mathbb{R}^{N},\\
U^{(s)} &\in
\mathbb{R}^{n_{\text{trials}} \times L \times N}.
\end{align}

For trial $i$, time $t$, and neuron $j$:

\begin{equation}
u_{i,t,j} = Z_i(t) \cdot \bm{b}_j + c_j.
\end{equation}

\subsection{Baseline}

The baseline vector is:

\begin{equation}
\bm{c}^{(s)} = (c_1^{(s)},\dots,c_N^{(s)}).
\end{equation}

The baseline sets the firing-rate offset before the nonlinearity. In the
poster simulation, all neurons use the same baseline drive. This keeps the
first experiment simple, but neuron-specific baselines are a natural extension
when matching real recordings.

\section{Firing Rate}

The neural drive $U$ can be negative, so it cannot be used directly as a firing
rate. A nonlinearity maps it into a positive rate:

\begin{equation}
\lambda_{i,t,j} = f(u_{i,t,j}).
\end{equation}

The poster simulation uses the softplus nonlinearity:

\begin{equation}
f(x) = \log(1 + e^x).
\end{equation}

The resulting tensor has shape:

\begin{equation}
\lambda^{(s)} \in
\mathbb{R}^{n_{\text{trials}} \times L \times N}.
\end{equation}

The rate $\lambda_{i,t,j}$ is measured in expected spikes per unit time for
neuron $j$ at time $t$ in trial $i$.

\section{Spike-Count Generation}

The base observation model samples spike counts as:

\begin{equation}
X_{i,t,j} \sim \text{Poisson}(\lambda_{i,t,j}\Delta t).
\end{equation}

Here $\lambda_{i,t,j}\Delta t$ is the expected spike count in one time bin.
Thus:

\begin{equation}
X^{(s)} \in
\mathbb{R}^{n_{\text{trials}} \times L \times N}.
\end{equation}

The pure Poisson model assumes that, conditional on the rate, spike counts
fluctuate randomly around their expected value. The temporal structure in $X$
comes from the temporal structure in $\lambda$, which itself comes from $Z$.

The simulator also includes optional phenomenological extensions that make the
observation model less idealized.

\subsection{Overdispersion}

Real spike counts often show variance larger than the mean. The pure Poisson
model cannot represent this because it implies:

\begin{equation}
\operatorname{Var}(X) = \mathbb{E}[X].
\end{equation}

To introduce overdispersion, the simulator can multiply the rate by a random
Gamma gain:

\begin{equation}
\lambda^{\ast}_{i,t,j} = \lambda_{i,t,j} g_{i,t,j},
\end{equation}

where:

\begin{equation}
\mathbb{E}[g_{i,t,j}] = 1.
\end{equation}

Spike counts are then sampled as:

\begin{equation}
X_{i,t,j} \sim \text{Poisson}(\lambda^{\ast}_{i,t,j}\Delta t).
\end{equation}

Setting the overdispersion parameter to zero recovers the pure Poisson
observation model.

\subsection{Refractory Period}

The simulator can approximate a refractory period at the bin level. After a
nonzero spike count, the following bins for the same neuron and trial can be
forced to zero for a sampled refractory duration.

The refractory duration is sampled as:

\begin{equation}
r_{\text{refractory}} \sim
\max\left(0, \operatorname{round}\left[
\mathcal{N}(\mu_r,\sigma_r^2)
\right]\right).
\end{equation}

This is not a biophysical membrane model. It is a statistical correction that
reduces very short inter-spike intervals.

\subsection{Bursts and Local Spike Clusters}

The simulator can also add short local spike clusters. A nonzero spike bin can
start a burst with probability:

\begin{equation}
p_{\text{burst}}.
\end{equation}

If a burst starts, extra spikes are sampled and added within a short temporal
window.

This models local spike clustering phenomenologically, without attempting to
simulate the cellular mechanisms that produce bursting.

\section{Windowing}

The spike tensor is reshaped:

\begin{equation}
X_{\text{reshaped}} \in
\mathbb{R}^{(n_{\text{trials}}L) \times N}.
\end{equation}

Temporal windows have length $W=10$ bins and are centered on each original
time bin.

With centered padding, one window is produced for every original time bin. The
missing samples at the beginning and end of each trial are filled with zeros.
Thus the windowed dataset preserves the full time-series length:

\begin{equation}
n_{\text{windows}} = n_{\text{trials}} L.
\end{equation}

This lets the encoder produce one embedding for each original time bin while
still receiving a local temporal context. For lag analysis, each window keeps
its trial identity, trial-local time index, and task condition. The condition
assigned to a window is the condition of the trial from which the window was
drawn.

Absolute trial-local time is used for lag-aware cross-population alignment.
Relative time in $[0,1]$ is appropriate for behavioral similarity, but
bin-level lag recovery requires integer trial-local time indices.

\section{Distance and Similarity Construction}

The circular reaching task uses two distances between windows:

\begin{itemize}
    \item temporal distance, based on  time
    \item circular condition distance, based on motor condition.
\end{itemize}

They are combined into a single task-informed distance. This means that two
windows are considered close if they are close in trial phase, close in motor
direction, or both.

More generally, the important distinction is not whether a variable is a
``single label'' or a ``multi-label'' target. The important distinction is the
geometry of each behavioral variable. A behavioral variable can be:

\begin{itemize}
    \item categorical, where equality or inequality is the relevant relation;
    \item circular, such as reaching direction on a circle;
    \item continuous one-dimensional, such as speed or position on a linear
    track;
    \item continuous multi-dimensional, such as a two-dimensional position
    $(x,y)$.
\end{itemize}

For a task containing both direction and two-dimensional position, the same
framework can be written as:

\begin{equation}
D_{\text{total}} =
w_tD_{\text{time}}
+ w_dD_{\text{circular}}(\text{direction})
+ w_pD_{\text{euclidean}}(\text{position}_x,\text{position}_y).
\end{equation}

The total distance is converted into a similarity matrix:

\begin{equation}
S = g(D_{\text{total}}).
\end{equation}

This similarity matrix can be used as supervision or structure for downstream
representation learning methods.

\section{Representation Learning Objective}

After windowing, each neural window is passed through an encoder:

\begin{equation}
z_i = f_\theta(X_i),
\end{equation}

where $X_i$ is a neural activity window and $z_i \in \mathbb{R}^3$ is the
learned representation. The poster comparison focuses on three models:

\begin{itemize}
    \item PCA: a linear unsupervised baseline that projects neural activity
    onto directions of maximal variance;
    \item temporal CNN1D: a local temporal encoder that uses convolutional
    kernels to detect short-range spike-pattern structure within each window;
    \item Transformer: an attention-based temporal encoder that can relate
    different bins within the window.
\end{itemize}

The input convention for the neural encoders is a batch of windows with shape:

\begin{equation}
\text{batch} \times \text{features} \times \text{sequence length}
=
\text{batch} \times N \times \text{window}.
\end{equation}

Thus the CNN is one-dimensional along time; neurons are treated as input
features/channels.

\subsection{Why These Three Encoders}

The three encoders play different scientific roles. They should not be read as
three arbitrary machine-learning models, but as three increasingly flexible
ways of asking the same question: how much of the shared latent task structure
can be recovered from population spike counts?

\subsubsection{PCA: Linear Variance-Based Baseline}

Principal Component Analysis is the simplest reference model. Given a matrix of
neural observations $X \in \mathbb{R}^{m \times N}$, where rows are observations
and columns are neurons, PCA finds orthogonal directions $v_1,\ldots,v_d$ that
maximize projected variance:

\begin{equation}
v_1 =
\arg\max_{\|v\|=1}
\operatorname{Var}(Xv).
\end{equation}

The low-dimensional representation is:

\begin{equation}
z_i = x_i V_d,
\end{equation}

where $V_d$ contains the first $d$ principal axes. In the poster experiment
$d=3$ when PCA is compared with the neural encoders.

PCA is useful because it is transparent. If PCA works well, then much of the
latent structure is already expressed as large-scale linear covariance in the
population activity. If PCA fails while nonlinear encoders work, then the
latent structure may be present but not linearly aligned with the dominant
variance directions.

Its limitation is equally important. PCA does not know about time, trial
structure, task condition, or lag. It only sees variance. A high-variance
direction may reflect nuisance variability rather than task geometry.
Therefore, PCA is a baseline, not the main model.

\subsubsection{Temporal CNN1D: Local Pattern Encoder}

The temporal CNN1D receives a neural window
$X_i \in \mathbb{R}^{N \times W}$, where $N$ is the number of neurons and $W$
is the window length in time bins. A one-dimensional convolution applies
temporal filters across the window:

\begin{equation}
h_{r,t}
=
\sigma
\left(
b_r +
\sum_{n=1}^{N}
\sum_{\Delta=0}^{K-1}
a_{r,n,\Delta}
X_{n,t+\Delta}
\right),
\end{equation}

where $r$ indexes filters, $K$ is the convolutional kernel size, and $\sigma$
is a nonlinearity. The key point is that the convolution runs along time. The
neurons are input features/channels.

This model is well matched to the spike-window setting. It can detect local
temporal patterns such as short increases in firing, local rate changes, and
coordinated patterns across neurons inside the window. It also has an
inductive bias: nearby time bins matter more directly than distant bins. This
is useful when the latent state evolves smoothly and when the window is meant
to capture local neural context around one time point.

In the current benchmark, the CNN1D is expected to be a strong compromise:
more expressive than PCA, but less flexible than a Transformer. If it recovers
a cleaner manifold and a plausible lag, that suggests the relevant information
is local in time and can be captured by short temporal filters.

The main limitation is window scale. If the window is too short, the CNN sees
too little temporal context. If the window is too long, the local filters may
mix irrelevant parts of the trial. Thus window length is not just a technical
hyperparameter; it is a modeling assumption about the temporal scale of neural
information.

\subsubsection{Transformer: Attention-Based Window Encoder}

The Transformer treats the bins inside a window as a small sequence. Each time
bin is mapped to a token representation, and self-attention computes how much
each bin should attend to the others:

\begin{equation}
\operatorname{Attention}(Q,K,V)
=
\operatorname{softmax}
\left(
\frac{QK^\top}{\sqrt{d_k}}
\right)V.
\end{equation}

Unlike a convolution, attention is not restricted to nearby positions. Within
the window, any bin can directly interact with any other bin. This makes the
Transformer more flexible: it can learn nonlocal relationships inside the
window and can, in principle, detect temporal patterns that are not captured by
fixed local kernels.

This flexibility is both a strength and a risk. A Transformer can align two
subjects well in terms of lag while producing an embedding geometry that is
compressed or visually less faithful to the original latent trajectory. In
that case, the model may be learning a representation that is useful for
cross-population consistency but not necessarily an interpretable recovery of
the ground-truth latent space.

For this reason, the Transformer should be evaluated with multiple diagnostics:
latent recovery, condition organization, Procrustes-aligned cross-population
similarity, and lag recovery. A good lag score alone is not sufficient.

\subsection{How To Read The Model Comparison}

The three models answer different diagnostic questions:

\begin{itemize}
    \item PCA asks whether the latent structure is visible through linear
    population covariance.
    \item CNN1D asks whether local temporal spike patterns are sufficient to
    recover task geometry and lag.
    \item Transformer asks whether a more flexible attention model can exploit
    within-window temporal relationships beyond local filters.
\end{itemize}

The best model is therefore not automatically the one with the best single
metric. A useful model should recover the known latent geometry, preserve
condition structure, and localize the imposed delay between the two
populations. These criteria can disagree, and that disagreement is itself
scientifically informative.

\subsection{Soft Structured Contrastive Cross-Entropy}

The main learning objective is a soft structured contrastive cross-entropy.
Its role is not to reconstruct $Z$ directly. Instead, task condition and
trial-time variables define a target distribution over which neural windows
should be close or far in the learned embedding.

For a batch of embeddings, the model computes cosine logits:

\begin{equation}
\ell_{ij} =
\frac{\operatorname{cos}(z_i,z_j)}{\tau_{\text{temp}}},
\end{equation}

where $\tau_{\text{temp}}$ is the contrastive temperature.

In a hard contrastive objective, a pair is usually either positive or negative.
Here the target relationship is soft. A pair can be strongly similar, weakly
similar, or dissimilar depending on task/time proximity. This is the same
softmax-denominator logic used in InfoNCE, but with a soft target distribution
instead of a single positive sample or a hard positive set.

The target similarity matrix $S$ is obtained from a metadata distance:

\begin{equation}
S_{ij} = \exp\left(-\frac{D_{\text{total},ij}}{\tau_S}\right).
\end{equation}

The total distance is a weighted combination of time and task-condition
distances:

\begin{equation}
D_{\text{total},ij}
=
w_t D^{\text{time}}_{ij}
+
w_c D^{\text{condition}}_{ij},
\qquad
w_t+w_c=1.
\end{equation}

The time distance is:

\begin{equation}
D^{\text{time}}_{ij} = |t_i-t_j|.
\end{equation}

The condition distance can be defined in different ways. For the current
circular motor task, the run uses circular condition distance:

\begin{equation}
D^{\text{condition}}_{ij}
=
\min\left(
|c_i-c_j|,
n_{\text{conditions}}-|c_i-c_j|
\right).
\end{equation}

This choice makes neighboring directions more similar than opposite
directions. If the intended supervision should only express same/different
condition information, the appropriate alternative is categorical distance:

\begin{equation}
D^{\text{condition}}_{ij}
=
\begin{cases}
0, & c_i=c_j,\\
1, & c_i\ne c_j.
\end{cases}
\end{equation}

This distinction is scientifically important. Circular distance injects a weak
task geometry over conditions. Categorical distance only imposes class consistency
and does not impose circular ordering.

For each anchor $i$, the row $S_{i,:}$ is normalized into a target probability
distribution over other batch elements. This target distribution is denoted
$q_{ij}$:

\begin{equation}
q_{ij} =
\frac{S_{ij}}{\sum_{k \ne i} S_{ik}},
\quad j \ne i.
\end{equation}

The embedding logits define the model probability $p_{ij}$, i.e. the
similarity distribution induced by the current learned embeddings:

\begin{equation}
p_{ij} =
\frac{\exp(\ell_{ij})}{\sum_{k \ne i}\exp(\ell_{ik})}.
\end{equation}

The loss is the cross-entropy between the soft target distribution $q_{ij}$
and the embedding-induced distribution $p_{ij}$. For a minibatch of size $B$:

\begin{equation}
\mathcal{L}_{\text{soft-CE}} =
- \frac{1}{B}
\sum_{i=1}^{B}
\sum_{\substack{j=1 \\ j \ne i}}^{B}
q_{ij}\log p_{ij}.
\end{equation}

Here $B$ is the number of neural windows in the minibatch. The matrix
$Q=(q_{ij})$ is the target similarity distribution induced by task/time
distances, while $P=(p_{ij})$ is the similarity distribution induced by the
current embeddings.

This objective is useful because it can use graded information: two windows can
be partially similar if they are close in time, close in task condition, or
both. It sits between a purely temporal InfoNCE objective and a hard supervised
contrastive objective. Its main risk is that the learned geometry may reflect
the chosen metadata weights too strongly. If the time term dominates, the
embedding may mostly encode trial phase. If the condition term dominates, the
embedding may separate conditions without preserving continuous trajectory
geometry.

Therefore, this loss is task-structured contrastive learning. It is not a claim
of purely unsupervised discovery: the simulator is used to test whether the
representation learned under these constraints also recovers the known latent
geometry and the known cross-population lag.

\subsection{Supervised InfoNCE Baseline}

The framework also supports a more standard supervised InfoNCE objective. In
this case, positives are all samples in the batch with the same condition
label:

\begin{equation}
\mathcal{P}(i) =
\{j \ne i : y_j = y_i\}.
\end{equation}

The loss is:

\begin{equation}
\mathcal{L}_{\text{supNCE}} =
- \frac{1}{B}\sum_i
\frac{1}{|\mathcal{P}(i)|}
\sum_{p \in \mathcal{P}(i)}
\log
\frac{\exp(\operatorname{cos}(z_i,z_p)/\tau)}
{\sum_{a \ne i}\exp(\operatorname{cos}(z_i,z_a)/\tau)}.
\end{equation}

This objective is easier to communicate and closer to standard supervised
contrastive learning. However, it treats all same-label samples as equally
positive and ignores the circular distance between neighboring directions and
the temporal structure of the trial.

\subsection{Unsupervised Time-Offset InfoNCE}

The framework also supports an unsupervised temporal baseline. In this
objective, behavioral labels are not used. A positive pair is defined by
temporal offset within the same trial:

\begin{equation}
\mathcal{P}_{\Delta}(i) =
\{j : \operatorname{trial}(j)=\operatorname{trial}(i),
|t_j-t_i|=\Delta\}.
\end{equation}

The same InfoNCE denominator is then used over the batch. This is important as
a baseline because it asks whether the network can recover a useful state
space without explicit task labels.

\section{What the Benchmark Tests}

The benchmark separates three objects that are usually entangled in real data:
the task-level latent geometry, the population-specific neural readout, and the
stochastic spike-count observation. This separation is useful because the true
latent trajectory is known. We can therefore ask whether a learned embedding
recovers the latent geometry, whether it preserves condition structure, and
whether it reveals the imposed temporal offset between the two populations.

The scientific question is not whether an embedding looks visually clean. The
question is whether a model recovers information that was present in the
controlled generative process:

\begin{itemize}
    \item latent geometry: does the embedding preserve the low-dimensional task
    trajectory?
    \item condition structure: do trajectories remain organized by movement
    direction?
    \item temporal structure: does the embedding preserve within-trial phase?
    \item cross-subject structure: can two subject-specific neural populations
    be aligned back to a common task space?
    \item lag structure: can the known temporal delay between populations be
    recovered from the learned manifolds?
\end{itemize}

This distinction matters because a model can recover one part of the structure
and fail on another. For example, it may localize the correct inter-subject lag
while still producing a compressed or distorted latent geometry.

\section{Multi-Subject Shared Latent Setting}

The simulator can generate multiple subject-specific neural observations from a
shared latent/task process. The screen stimulus, trial labels, and task-level
latent trajectory are common to both subjects:

\begin{equation}
Z_{\text{task}}(t)
\rightarrow
X^{(1)}(t), X^{(2)}(t), \dots
\end{equation}

Each subject has a different loading matrix, baseline, and spike observation
noise:

\begin{equation}
U^{(s)} = Z^{(s)} B^{(s)} + \bm{c}^{(s)}.
\end{equation}

The task labels and motor state are shared, but the neural response can be
subject-specific. It is useful to distinguish two latent objects:

\begin{itemize}
    \item $Z_{\text{task}}$: the shared task/stimulus latent trajectory used as
    the ground-truth target for evaluation;
    \item $Z_{\text{neural-driver}}^{(s)}$: the subject-specific latent driver,
    possibly lagged, used to generate neural activity.
\end{itemize}

This creates a controlled test of whether representation learning can recover a
common latent/task geometry from different neural observations.

\subsection{Temporal Lag}

In interactive motor settings, two subjects may share the same task structure
but show a small temporal offset in neural activity. For example, one subject
may act while another observes or responds with delay.

The simulator represents this using a subject-specific neural response lag:

\begin{equation}
Z_{\text{neural-driver}}^{(s)}(t) = Z_{\text{task}}(t - \ell_s),
\end{equation}

where $\ell_s$ is measured in time bins.

The labels remain unchanged because both subjects see the same screen and task
events. Only the neural driver used to generate spikes is shifted. This makes it
possible to test whether learned embeddings remain alignable when
subject-specific neural observations are slightly shifted in time.

\section{Shared-Space Alignment}

If two encoders recover the same latent structure, their coordinates need not
match exactly. Latent representations are often identifiable only up to
geometric transformations such as translation, rotation, reflection, scaling,
axis permutation, mild nonlinear warping, and temporal lag.

Therefore, cross-subject evaluation should not require raw coordinate equality.
Instead, it should test whether the recovered geometries are equivalent after
reasonable alignment.

\subsection{Recommended Evaluation}

A conservative evaluation stack is:

\begin{itemize}
    \item center and scale embeddings;
    \item align embeddings using orthogonal Procrustes;
    \item measure residual alignment error;
    \item compare representational distance matrices;
    \item compute latent recovery against the known $Z$;
    \item search over small temporal lags when the task implies delayed
    interaction.
\end{itemize}

More flexible methods, such as affine alignment, CCA, CKA, dynamic time
warping, manifold alignment, optimal transport, or Gromov--Wasserstein
alignment, can be useful later. However, overly flexible alignment can hide real
failures of the representation model. The first benchmark should therefore
start with Procrustes alignment and representational similarity analysis.

\subsection{Lag-Aware Procrustes Alignment}

For cross-subject comparison, the current baseline estimates the temporal
offset by scanning a small set of candidate lags:

\begin{equation}
\ell \in \{-15,\dots,0,\dots,15\}.
\end{equation}

For each candidate lag, the comparison pairs:

\begin{equation}
z^{(1)}(\text{trial},t)
\quad \text{with} \quad
z^{(2)}(\text{trial},t+\ell).
\end{equation}

The second subject embedding is then aligned to the first subject embedding
using orthogonal Procrustes, and the lag with the best Procrustes $R^2$ is
selected:

\begin{equation}
\ell^\ast =
\arg\max_{\ell} R^2_{\text{Procrustes}}(\ell).
\end{equation}

In the current synthetic configuration, the delayed population is generated
with a ten-bin neural-response lag:

\begin{equation}
\ell_{\text{delayed}} = 10 \quad \text{bins} = 100 \quad \text{ms}.
\end{equation}

The CNN gives the most sensible trade-off between latent geometry recovery and
lag localization. The Transformer can peak at the correct lag while producing a
compressed and less faithful geometry. This dissociation is important:
cross-subject temporal consistency and ground-truth latent recovery are not the
same validation target.

The estimated lag is selected over candidate lags, not assumed.


\subsection{Poster Pipeline: From Spikes to Temporal Asymmetry}

The complete poster pipeline can be written as:

\begin{align}
Z_{\text{task}}
&\rightarrow
\left(
Z_{\text{driver}}^{(\text{earlier})},
Z_{\text{driver}}^{(\text{delayed})}
\right)\\
&\rightarrow
\left(
U^{(\text{earlier})}, U^{(\text{delayed})}
\right)
\rightarrow
\left(
\lambda^{(\text{earlier})}, \lambda^{(\text{delayed})}
\right)\\
&\rightarrow
\left(
X^{(\text{earlier})}, X^{(\text{delayed})}
\right)
\rightarrow
\text{windows}
\rightarrow
\left(
E^{(\text{earlier})}, E^{(\text{delayed})}
\right)\\
&\rightarrow
\text{lag diagnostics}.
\end{align}

The shared latent is:

\begin{equation}
Z_{\text{task}} \in \mathbb{R}^{200 \times 100 \times 3}.
\end{equation}

The subject-specific neural drivers are:

\begin{align}
Z_{\text{driver}}^{(\text{earlier})}(t) &= Z_{\text{task}}(t),\\
Z_{\text{driver}}^{(\text{delayed})}(t) &= Z_{\text{task}}(t-10).
\end{align}

Thus the delayed population is shifted by $10$ bins, corresponding to
$100$ ms.

The spike tensors are:

\begin{align}
X^{(\text{earlier})} &\in \mathbb{R}^{200 \times 100 \times 80},\\
X^{(\text{delayed})} &\in \mathbb{R}^{200 \times 100 \times 80}.
\end{align}

The encoders operate on center-padded neural windows:

\begin{equation}
W_i^{(s)} \in \mathbb{R}^{10 \times 80}.
\end{equation}

Each window is mapped to a three-dimensional embedding:

\begin{equation}
E_i^{(s)} = f_\theta(W_i^{(s)}),
\qquad
E_i^{(s)} \in \mathbb{R}^{3}.
\end{equation}

Because the stride is one and padding preserves trial length, each subject has:

\begin{equation}
200 \times 100 = 20000
\end{equation}

window embeddings.

\subsection{Manifold-Level Temporal Diagnostic}

Beyond the Procrustes lag scan, the recovered manifolds can be inspected with a
lag-lag diagnostic. The goal is to ask whether temporal structure in the
earlier population appears shifted relative to temporal structure in the
delayed population.

For each population, the recovered embedding sequence is divided into temporal
blocks. For block $b$, let:

\begin{equation}
M_b^{(\text{earlier})}, M_b^{(\text{delayed})}
\in \mathbb{R}^{n_{\text{obs}} \times d_{\text{emb}}},
\end{equation}

where rows are trial/block observations and columns are embedding dimensions.
A simple diagnostic compares block summaries across populations:

\begin{equation}
L_{b,b'} =
\operatorname{sim}
\left(
M_b^{(\text{earlier})},
M_{b'}^{(\text{delayed})}
\right).
\end{equation}

In the current exploratory version, the similarity can be a correlation between
block-level summaries. This should be interpreted cautiously: it is a visual
diagnostic of temporal asymmetry, not a causal coupling measure.

\section{Poster Storyline}

This section explains how to read and present the poster. The goal is not only
to describe the figures, but to make clear what scientific claim each block is
allowed to support.

\subsection{Question}

The poster starts from one question: if two neural populations are driven by
the same task-level latent process, but through different neural readouts and
with a small temporal delay, can representation learning recover a shared
geometry and reveal the delay?

The important point is that the shared latent space is known because the data
are synthetic. This gives a controlled benchmark. We are not simply asking
whether an embedding looks organized. We are asking whether it recovers a
known ground-truth structure and whether temporal asymmetry can be detected
after encoding.

\subsection{Workflow}

The workflow summarizes the full causal chain used in the experiment:

\begin{align}
\text{shared task latent}
&\rightarrow
\text{subject-specific readouts}
\rightarrow
\text{spike populations}
\rightarrow
\text{windows}
\\
&\rightarrow
\text{encoders}
\rightarrow
\text{recovered manifolds and lag diagnostics}.
\end{align}

This order matters. The two populations do not start from identical spike
matrices. They share the same task-level latent process, but each population
has its own neural loading matrix and stochastic spike generation. Therefore,
any recovered similarity between the populations must survive the
subject-specific neural readout and the spike-count observation process.

\subsection{Generative Block}

The generative block explains how the synthetic neural data are created. The
shared latent trajectory is denoted by $Z$. For each population, $Z$ is mapped
to a neural drive:

\begin{equation}
U = ZB + \bm{c},
\end{equation}

where $B$ is a population-specific loading matrix and $\bm{c}$ is an offset.
The drive is transformed into a positive firing intensity through a softplus
nonlinearity:

\begin{equation}
\lambda = \operatorname{softplus}(U).
\end{equation}

Spike counts are then sampled from a Poisson observation model:

\begin{equation}
X \sim \operatorname{Poisson}(\lambda \Delta t),
\end{equation}

with optional phenomenological corrections such as overdispersion,
refractoriness, and short burst-like clusters. The delayed population is
generated from a temporally shifted neural driver. In the poster run, the
delay is ten bins, corresponding to $100$ ms.

\subsection{Panels A and B}

Panel A shows the ground-truth latent trajectories. These are the trajectories
that the learning models should recover from spikes. The colors correspond to
motor directions. Because $Z$ is known, the figure is not an exploratory
visualization: it is the target geometry of the benchmark.

Panel B shows observed spike rasters from the two generated populations. This
panel is deliberately much less structured visually than Panel A. It reminds
the reader that the encoder does not receive the latent trajectory directly.
It receives noisy, binned spike counts generated through population-specific
readouts.

\subsection{Learning Block}

The learning block is the methodological center of the poster. Each spike
matrix is converted into local temporal windows. In the poster configuration,
each window contains ten time bins around a target bin and all recorded neurons.
The encoder maps each window to a low-dimensional embedding:

\begin{equation}
e_i = f_\theta(W_i),
\qquad
e_i \in \mathbb{R}^{3}.
\end{equation}

Three encoder families are compared:

\begin{itemize}
    \item PCA, which asks whether the latent geometry is already visible in
    linear population covariance;
    \item CNN1D, which asks whether local temporal spike patterns are
    sufficient to recover geometry and delay;
    \item Transformer, which asks whether attention over the window can exploit
    more flexible temporal relationships.
\end{itemize}

The contrastive target is not a hard same/different label. Instead, a
metadata distance is built from trial time and task condition:

\begin{equation}
D_{ij}
=
w_t D^{\text{time}}_{ij}
+
w_c D^{\text{condition}}_{ij}.
\end{equation}

This distance is converted into a soft target distribution:

\begin{equation}
q_{ij}
\propto
\exp(-D_{ij}/\tau_S).
\end{equation}

The embedding itself induces a predicted similarity distribution:

\begin{equation}
p_{ij}
\propto
\exp
\left(
\operatorname{cos}(e_i,e_j)/\tau
\right).
\end{equation}

Training minimizes cross-entropy between the target distribution $Q$ and the
embedding-induced distribution $P$:

\begin{equation}
\mathcal{L}_{\text{soft-CE}}
=
-
\frac{1}{B}
\sum_{i=1}^{B}
\sum_{\substack{j=1 \\ j \ne i}}^{B}
q_{ij}\log p_{ij}.
\end{equation}

This is a task-structured contrastive objective. It is related to InfoNCE
because it uses a softmax distribution over other samples in the batch, but it
uses a soft target distribution rather than a single positive sample.

\subsection{Panels C and D}

Panels C and D show recovered manifolds from the CNN encoder. Panel C is shown
in three dimensions because the encoder output has dimension three and the
goal is to inspect the learned geometry in its native output space. Panel D is
shown as a two-dimensional projection for readability: it gives a cleaner view
of condition organization for the delayed population.

The difference between a 3D view and a 2D projection should not be
overinterpreted. A 2D projection is a visualization choice. The model output
can still be three-dimensional even when the plotted figure displays only two
axes or a projection. The key question is whether the recovered trajectories
preserve condition structure and resemble the known latent geometry.

\subsection{Temporal Recovery}

Temporal recovery asks whether the delay imposed during generation can be
estimated from the learned embeddings. The earlier and delayed populations are
not compared point-by-point without adjustment. Instead, the delayed embedding
is shifted by a set of candidate lags:

\begin{equation}
\ell \in \{-15,\ldots,15\}.
\end{equation}

For each candidate lag, the delayed embedding is paired with the earlier
embedding at the corresponding shifted time. The two embeddings are then
aligned with orthogonal Procrustes.

\subsection{Procrustes Alignment}

Procrustes alignment answers a geometric question: are two recovered spaces
the same up to a rigid transformation? Given two centered matrices $A$ and
$B$, Procrustes finds the orthogonal matrix $R$ that minimizes:

\begin{equation}
\|A - BR\|_F^2
\quad
\text{subject to}
\quad
R^\top R = I.
\end{equation}

This allows rotation and reflection, but does not allow arbitrary warping.
That is why it is useful as a conservative alignment diagnostic. If two
embeddings align well after Procrustes, they share a similar geometry even if
their axes are not identical.

In the lag scan, this alignment is repeated for every candidate lag. The
selected lag is the lag with the best Procrustes alignment score:

\begin{equation}
\ell^\ast
=
\arg\max_{\ell}
R^2_{\text{Procrustes}}(\ell).
\end{equation}

\subsection{Panel E: Lag Recovery Scan}

Panel E shows the alignment score as a function of candidate lag. The vertical
reference line marks the true imposed delay of ten bins. A model that peaks
near this value has recovered the temporal offset between the two populations.

However, this panel should be interpreted together with the manifold plots. A
model can recover the correct lag while still producing a distorted or
compressed geometry. Therefore, lag recovery is necessary but not sufficient
for claiming faithful latent recovery.

\subsection{Lag-Lag Map}

The lag-lag map is a complementary temporal diagnostic. Instead of selecting a
single candidate delay, both recovered manifolds are divided into temporal
blocks. For each earlier block $b$ and delayed block $b'$, a similarity score
is computed:

\begin{equation}
L_{b,b'}
=
\operatorname{sim}
\left(
M_b^{(\text{earlier})},
M_{b'}^{(\text{delayed})}
\right).
\end{equation}

This produces a matrix of temporal correspondences. If the two populations
were perfectly synchronous, the strongest structure would be expected near the
main diagonal. If one population is delayed, the structure may shift away from
the diagonal. In the current poster this is exploratory: it is meant to
visualize temporal asymmetry, not to prove causal coupling.

\subsection{Panel F: Exploratory Lag-Lag Map}

Panel F shows this block-wise temporal comparison. The map should be read as a
summary of how temporal profiles in one recovered manifold relate to temporal
profiles in the other. Warm or cool regions indicate positive or negative
correspondence between blocks, depending on the chosen similarity measure.

The main interpretive point is modest: the lag-lag map can reveal temporal
structure that is not visible in a single static embedding plot. It is not yet
a decoder, not a causal analysis, and not a replacement for cross-validated
prediction.

\subsection{Interpretation and Required Controls}

The current interpretation is that the benchmark can separate three questions:

\begin{itemize}
    \item whether the latent task geometry is recoverable from spikes;
    \item whether the two population embeddings can be aligned into a common
    space;
    \item whether the imposed delay can be estimated from the recovered
    manifolds.
\end{itemize}

Before making stronger claims, several controls are required:

\begin{itemize}
    \item a lag-zero control, where both populations are generated without
    temporal delay;
    \item repeated simulations across multiple random seeds;
    \item time-only and condition-only loss ablations;
    \item shuffled-condition and shuffled-time controls;
    \item comparison between PCA, CNN1D, and Transformer under identical
    training and evaluation settings;
    \item cross-validated prediction from one population embedding to the
    other population, which would be stronger than a visual lag-lag map.
\end{itemize}

These controls are essential because a visually appealing manifold is not
enough. The goal is to show that the recovered geometry and temporal asymmetry
are stable consequences of the generative structure and learning objective,
not artifacts of plotting, seed choice, or overly strong metadata supervision.

\section{Operational Questions For The Poster Discussion}

This section collects the questions that are most likely to arise during a
poster discussion. The emphasis is practical: what object is manipulated, what
the model receives, and what each plot is allowed to mean.

\subsection{What Does It Mean To Lag One Population?}

In the poster simulation, the delayed population is not delayed by
$10$ milliseconds. It is delayed by $10$ time bins. Since the bin size is
$\Delta t = 10$ ms, this corresponds to:

\begin{equation}
10 \text{ bins} \times 10 \text{ ms/bin} = 100 \text{ ms}.
\end{equation}

The lag is applied at the level of the latent neural driver, before spike
generation. The shared task trajectory is $Z_{\text{task}}(t)$. The earlier
population uses:

\begin{equation}
Z_{\text{driver}}^{(\text{earlier})}(t)
=
Z_{\text{task}}(t),
\end{equation}

whereas the delayed population uses:

\begin{equation}
Z_{\text{driver}}^{(\text{delayed})}(t)
=
Z_{\text{task}}(t-10).
\end{equation}

Intuitively, the delayed population is driven by the same latent trajectory,
but it is seeing an earlier state of that trajectory at the current time bin.
If the earlier population is at latent state $Z(t)$, the delayed population is
approximately at the state that the earlier population had ten bins before.

At the beginning of the trial, there is no time $t-10$. The implementation does
not wrap the trajectory around the trial. Instead, it pads the beginning using
the edge value of the trial. In code, the positive-lag case is equivalent to:

\begin{align}
Z_{\text{lagged}}[:, :\ell, :]
&=
Z[:, :1, :],
\\
Z_{\text{lagged}}[:, \ell:, :]
&=
Z[:, :-\ell, :].
\end{align}

Thus the lag is within trial, not across trials. Trial labels and task
conditions remain the same for both populations, because both populations are
assumed to share the same external task events. What changes is the neural
response timing.

\subsubsection{Matrix View Of The Lag}

For one trial, ignore the trial index and write the latent trajectory as a
matrix:

\begin{equation}
Z =
\begin{bmatrix}
z_1 \\
z_2 \\
\cdots \\
z_L
\end{bmatrix}
\in
\mathbb{R}^{L \times k},
\qquad
z_t \in \mathbb{R}^{k}.
\end{equation}

With a positive lag of $\ell$ bins, the delayed trajectory is:

\begin{equation}
Z^{(\text{delayed})}
=
\begin{bmatrix}
z_1 \\
\cdots \\
z_1 \\
z_1 \\
z_2 \\
\cdots \\
z_{L-\ell}
\end{bmatrix}
\in
\mathbb{R}^{L \times k},
\end{equation}

where the first $\ell$ rows are filled by the first valid latent state because
the current implementation uses edge padding. The important point is that the
shape does not change. The earlier and delayed populations still have latent
tensors with shape:

\begin{equation}
n_{\text{trials}} \times L \times k.
\end{equation}

The delayed population is therefore not a different task and not a different
label sequence. It is the same task-related latent trajectory sampled with a
response delay before generating spikes.

\subsection{What Is One Observational Unit?}

The raw spike tensor for one population has shape:

\begin{equation}
X \in \mathbb{R}^{n_{\text{trials}} \times L \times N},
\end{equation}

where $L$ is the number of time bins per trial and $N$ is the number of
neurons. In the poster configuration:

\begin{equation}
n_{\text{trials}}=200,
\quad
L=100,
\quad
N=80.
\end{equation}

The learning models do not receive an entire trial as one sample. They receive
local windows. With window size $W=10$, one observational unit is:

\begin{equation}
W_i \in \mathbb{R}^{10 \times 80}.
\end{equation}

This means: ten consecutive time bins, all neurons. With centered padding and
stride one, every original time bin has one associated window. Therefore one
population yields:

\begin{equation}
200 \times 100 = 20000
\end{equation}

windows and therefore 20000 embeddings.

\subsection{What Is A Batch?}

A batch is a collection of windows sampled for one optimization step. If the
batch size is $B$, the model receives:

\begin{equation}
\text{batch}
\in
\mathbb{R}^{B \times 10 \times 80}.
\end{equation}

So, if $B=30$, then the batch contains 30 windows, and each window contains 10
time bins by 80 neurons. In the current baseline suite the batch size is a
hyperparameter; for example, one configuration uses $B=256$. The important
point is:

\begin{equation}
\text{batch size} \ne \text{window size}.
\end{equation}

The window size says how much local temporal context each sample contains. The
batch size says how many such samples are processed together and compared
inside the contrastive loss.

\subsection{How Does Each Model See The Same Window?}

All models receive windows derived from the same spike data, but they interpret
the window differently.

\subsubsection{PCA}

PCA is not temporal in itself. To apply PCA to a window, the window is treated
as a vector of features. A $10 \times 80$ window can be flattened into:

\begin{equation}
\operatorname{vec}(W_i) \in \mathbb{R}^{800}.
\end{equation}

PCA then finds linear directions of maximal variance across all windows. It
does not know that the first ten groups of features came from ordered time
bins unless this temporal structure is already reflected in the covariance.

Therefore, PCA is useful as a baseline: it asks whether the latent structure is
already visible in linear variance. But PCA does not learn temporal filters,
attention, or task-aware nonlinear transformations.

\subsubsection{CNN1D}

The CNN1D treats the window as a short temporal signal. Conceptually, the
input is arranged as:

\begin{equation}
\text{batch} \times \text{neurons} \times \text{time}
=
B \times 80 \times 10.
\end{equation}

The convolution slides filters along the time dimension. A filter can detect
local temporal motifs such as a short rise in firing rate, a local pattern
across neurons, or a brief coordinated event. This is why the window is most
natural for the CNN: the model is explicitly designed to exploit local
temporal structure.

The CNN has a strong inductive bias. It assumes that nearby bins contain useful
information and that local temporal patterns can be reused across the window.
This can be an advantage when neural dynamics are smooth or locally structured.
It can be a limitation if the relevant relationship is nonlocal or depends on
longer trial context.

\subsubsection{Transformer}

The Transformer also receives the same $10 \times 80$ window, but interprets
it as a sequence of ten tokens. Each token corresponds to one time bin, whose
features are the neural population activity at that bin:

\begin{equation}
W_i =
(x_{t-4}, x_{t-3}, \ldots, x_{t+5}),
\qquad
x_t \in \mathbb{R}^{80}.
\end{equation}

Self-attention lets each time bin compare itself with every other bin inside
the window:

\begin{equation}
\operatorname{Attention}(Q,K,V)
=
\operatorname{softmax}
\left(
\frac{QK^\top}{\sqrt{d_k}}
\right)V.
\end{equation}

This means the Transformer is less local than the CNN. It can learn that the
beginning and end of the window are related, or that some bins should be
weighted more than others. This flexibility can help with temporal alignment,
but it can also produce embeddings that are useful for the lag score while
being geometrically compressed or harder to interpret.

\subsubsection{What A Token Means Here}

In language models, a token is usually a word fragment. In this neural setting,
a token is not a word. A token is one time bin of neural population activity.
For a window with $W=10$ bins and $N=80$ neurons, the input window is:

\begin{equation}
W_i =
\begin{bmatrix}
x_{t-4} \\
x_{t-3} \\
\cdots \\
x_{t+5}
\end{bmatrix},
\qquad
x_t \in \mathbb{R}^{80}.
\end{equation}

Thus the Transformer sees ten tokens:

\begin{equation}
\underbrace{x_{t-4}}_{\text{token 1}},
\underbrace{x_{t-3}}_{\text{token 2}},
\ldots,
\underbrace{x_{t+5}}_{\text{token 10}}.
\end{equation}

Each token is an 80-dimensional vector containing the activity of the whole
population at one time bin. A linear layer first maps each token from
$\mathbb{R}^{80}$ to the internal Transformer dimension $d_{\text{model}}$.
Self-attention then lets the ten bins inside the same window compare with one
another. After attention, the window is summarized into one output embedding:

\begin{equation}
f_\theta(W_i) = e_i,
\qquad
e_i \in \mathbb{R}^{3}.
\end{equation}

This is the crucial dimensional point: the Transformer does not output one
embedding per token in the final representation used for the contrastive loss.
It uses the ten tokens to compute one representation for the whole local
window.

\subsubsection{Do Windows Talk To Each Other Inside The Transformer?}

In the current setup, no. The Transformer operates inside one window at a
time. If the batch has shape:

\begin{equation}
B \times 10 \times 80,
\end{equation}

then the Transformer processes each of the $B$ windows independently and
returns:

\begin{equation}
E_{\text{batch}} \in \mathbb{R}^{B \times 3}.
\end{equation}

The communication between different windows happens in the loss, not inside
the Transformer. This gives two separate levels:

\begin{itemize}
    \item within-window learning: CNN filters or Transformer attention extract
    a representation from the ten-bin neural context;
    \item between-window learning: the contrastive loss compares the resulting
    embeddings across windows in the batch.
\end{itemize}

So the model first converts each local neural movie
$W_i \in \mathbb{R}^{10 \times 80}$ into one point
$e_i \in \mathbb{R}^{3}$. Then the loss asks whether the distances among those
points are compatible with the task/time metadata distances among the
corresponding windows.

\subsection{Why Use Windows For PCA And Transformer Too?}

Using windows for all models keeps the comparison fair: each method receives
the same local neural information around each target time bin. The difference
is not the data each method sees, but how the method processes that data.

For PCA, the window is a high-dimensional vector. For CNN1D, the window is a
local temporal signal. For the Transformer, the window is a short sequence.
Thus the same window supports three different modeling assumptions:

\begin{itemize}
    \item linear covariance structure for PCA;
    \item local temporal filters for CNN1D;
    \item flexible within-window temporal relations for the Transformer.
\end{itemize}

\subsection{What Exactly Does Procrustes Do In The Lag Plot?}

The Procrustes step is not used to train the model. It is used after training
as an evaluation tool. Suppose we have two embeddings:

\begin{equation}
E^{(\text{earlier})}
\quad \text{and} \quad
E^{(\text{delayed})}.
\end{equation}

Even if both embeddings recover the same latent geometry, their coordinate
axes may be rotated or reflected relative to each other. Procrustes asks:
can one embedding be rotated/reflected so that it overlaps the other?

For a candidate lag $\ell$, we first pair points by trial and shifted time:

\begin{equation}
E^{(\text{earlier})}(\text{trial},t)
\leftrightarrow
E^{(\text{delayed})}(\text{trial},t+\ell).
\end{equation}

Then we find the best orthogonal transformation $R_\ell$:

\begin{equation}
R_\ell
=
\arg\min_{R^\top R=I}
\left\|
E^{(\text{earlier})}
-
E^{(\text{delayed})}R
\right\|_F^2.
\end{equation}

The score plotted in the lag recovery graph is high when the delayed embedding,
after the candidate shift and Procrustes rotation, matches the earlier
embedding well. The best lag is:

\begin{equation}
\ell^\ast =
\arg\max_\ell
R^2_{\text{Procrustes}}(\ell).
\end{equation}

Intuitively, the graph says: ``if I shift the delayed population by this many
bins, how well can its recovered manifold be rotated onto the earlier
population's manifold?'' A peak near the true lag suggests that the embedding
contains the temporal offset. It does not by itself prove that the entire
latent geometry was recovered faithfully.

\subsubsection{What Is Lost When Testing A Candidate Lag?}

Suppose the two recovered embeddings are:

\begin{equation}
E^{(\text{earlier})},
E^{(\text{delayed})}
\in
\mathbb{R}^{20000 \times 3}.
\end{equation}

The number 20000 is the number of window-centered observations, for example
$200$ trials times $100$ bins. The number 3 is the embedding dimension.

If we test a candidate lag of $\ell=15$ bins, we cannot compare all 20000 rows,
because the shifted sequence has missing partners at the edges of each trial.
For a positive lag, the comparison is conceptually:

\begin{equation}
E^{(\text{earlier})}(\text{trial},t)
\quad \text{with} \quad
E^{(\text{delayed})}(\text{trial},t+\ell).
\end{equation}

Only valid within-trial pairs are kept. Therefore, some rows near the beginning
or end of each trial are discarded for that candidate lag. What is discarded is
time points, not embedding dimensions. The matrices passed to Procrustes are:

\begin{equation}
A_\ell, B_\ell \in \mathbb{R}^{m_\ell \times 3},
\end{equation}

where $m_\ell$ is the number of valid matched rows after shifting. The three
embedding coordinates are still present.

\subsubsection{What Procrustes Maximizes}

For each candidate lag $\ell$, Procrustes finds the best rotation/reflection
of the delayed embedding onto the earlier embedding:

\begin{equation}
R_\ell
=
\arg\min_{R^\top R=I}
\left\|
A_\ell - B_\ell R
\right\|_F^2.
\end{equation}

The associated score can be expressed as an alignment $R^2$: how much of the
variance in $A_\ell$ is explained by the optimally aligned $B_\ell$. The lag
scan repeats this for many candidate lags and selects:

\begin{equation}
\ell^\ast =
\arg\max_\ell R^2(\ell).
\end{equation}

Thus the lag score is not a nonlinear decoder. It is a lag-conditioned linear
geometric alignment test. It asks whether, after matching the correct temporal
rows and allowing a rigid change of coordinate system, the two recovered
manifolds occupy the same geometry.

\subsection{How To Explain The Soft Contrastive Loss}

The loss builds a target notion of similarity between windows before looking
at the embedding. Two windows are treated as similar if they are close in trial
time, close in task condition, or both. This is encoded by a distance:

\begin{equation}
D_{ij}
=
w_t D^{\text{time}}_{ij}
+
w_c D^{\text{condition}}_{ij}.
\end{equation}

This distance is transformed into a soft probability distribution:

\begin{equation}
q_{ij}
=
\frac{\exp(-D_{ij}/\tau_S)}
{\sum_{k\ne i}\exp(-D_{ik}/\tau_S)}.
\end{equation}

The embedding produces its own probability distribution:

\begin{equation}
p_{ij}
=
\frac{\exp(\operatorname{cos}(e_i,e_j)/\tau)}
{\sum_{k\ne i}\exp(\operatorname{cos}(e_i,e_k)/\tau)}.
\end{equation}

The loss minimizes:

\begin{equation}
\mathcal{L}_{\text{soft-CE}}
=
-
\frac{1}{B}
\sum_{i=1}^{B}
\sum_{j\ne i}
q_{ij}\log p_{ij}.
\end{equation}

In words: if the metadata say that two windows should be close, the embedding
is encouraged to assign them high similarity. If the metadata say they should
be far, the embedding is encouraged to assign them low similarity. This is
``soft'' because similarity is graded, not binary. It is contrastive because
each window is compared against other windows in the batch.

\subsubsection{Anchors, Soft Positives, And Effective Negatives}

In a standard hard contrastive setup, one often chooses an anchor, one positive
sample, and many negative samples. Here the logic is softer.

For a batch of $B$ windows, the encoder returns:

\begin{equation}
E_{\text{batch}}
=
\begin{bmatrix}
e_1 \\
e_2 \\
\cdots \\
e_B
\end{bmatrix}
\in
\mathbb{R}^{B \times 3}.
\end{equation}

For each row $i$, the embedding $e_i$ is the anchor. Every other row $j \ne i$
is a comparison candidate. The metadata distance $D_{ij}$ decides how much
candidate $j$ should behave like a positive example for anchor $i$:

\begin{itemize}
    \item small $D_{ij}$: window $j$ is a strong soft positive for window $i$;
    \item intermediate $D_{ij}$: window $j$ is partially related to window $i$;
    \item large $D_{ij}$: window $j$ behaves like an effective negative.
\end{itemize}

Thus the loss does not say ``same condition is positive, different condition
is negative'' in a hard binary way. It says that the whole similarity profile
around anchor $i$ should match the metadata-defined similarity profile around
that same anchor. This is why the loss can combine time and condition without
collapsing everything into a single hard label.

\subsubsection{Where The Window Ends And The Embedding Begins}

The model never compares raw windows directly in the loss. The comparison is
between embeddings. The complete pipeline for one batch is:

\begin{equation}
W_i \in \mathbb{R}^{10 \times 80}
\xrightarrow{\ f_\theta\ }
e_i \in \mathbb{R}^{3}
\xrightarrow{\ \text{batch comparison}\ }
\{p_{ij}\}_{j\ne i}.
\end{equation}

So, yes: the encoder first studies structure inside each window. Then the
loss studies relations between windows by comparing their 3D embeddings. The
metadata distance is defined between windows because each embedding
corresponds to one window.

\subsection{How To Explain The Lag-Lag Map}

The lag recovery scan searches for a single best delay. The lag-lag map asks a
more exploratory question: how do temporal blocks in one recovered manifold
relate to temporal blocks in the other?

Each recovered embedding sequence is divided into temporal blocks. For each
pair of blocks, one from the earlier population and one from the delayed
population, the analysis computes a similarity:

\begin{equation}
L_{b,b'}
=
\operatorname{sim}
\left(
M_b^{(\text{earlier})},
M_{b'}^{(\text{delayed})}
\right).
\end{equation}

Rows and columns are temporal blocks or lags. A diagonal pattern suggests that
similar temporal structure occurs at similar times. A shifted pattern suggests
that one population's temporal structure is displaced relative to the other.

The limitation is important: in the current poster, this is a descriptive
diagnostic. It summarizes temporal correspondence between recovered manifolds,
but it is not yet a decoding proof, not a causal measure, and not a statistical
test of coupling.

\subsection{Limits To State Clearly}

The main limits to state in discussion are:

\begin{itemize}
    \item the data are synthetic, so the current claim is about benchmark
    behavior, not direct biological discovery;
    \item the delay is imposed in the latent neural driver, not estimated from
    real simultaneous recordings;
    \item the soft contrastive loss injects task/time structure through the
    metadata distance, so the result is task-structured learning, not pure
    unsupervised discovery;
    \item a correct lag peak does not guarantee faithful recovery of the full
    latent geometry;
    \item the lag-lag map is exploratory and should be strengthened with
    cross-validated prediction or decoding;
    \item window size, metadata weights, temperature, batch size, and random
    seed need systematic sensitivity analysis.
\end{itemize}

\section{Poster Experimental Configuration}

The poster simulation uses the following concrete setting:

\begin{center}
\begin{tabular}{ll}
\hline
Quantity & Value\\
\hline
trials & $200$\\
trial length & $100$ bins\\
bin size & $\Delta t = 0.01$ s\\
trial duration & $1$ s\\
conditions & $8$ circular directions\\
latent dimension & $k=3$\\
neurons per population & $N=80$\\
temporal persistence & $\phi=0.4$\\
latent noise scale & $\sigma=0.15$\\
directional loading scale & $3.0$\\
extra loading scale & $0.051$\\
nonlinearity & softplus\\
overdispersion & $0.25$\\
refractory duration & mean $2$ bins, sd $1$ bin\\
burst probability & $0.05$\\
burst window & $3$ bins\\
earlier population lag & $0$ bins\\
delayed population lag & $10$ bins $=100$ ms\\
window size & $10$ bins\\
stride & $1$ bin\\
padding & centered zero padding\\
embedding dimension & $3$\\
contrastive temperature & $0.2$\\
soft similarity scale & $0.5$\\
time/label weights & $0.5/0.5$\\
\hline
\end{tabular}
\end{center}

For downstream analysis, each method can be exported as a trial-level structure
with one element per trial. For each trial, the relevant quantities are:

\begin{itemize}
    \item the trial direction;
    \item the earlier population spike matrix, $80 \times 100$;
    \item the delayed population spike matrix, $80 \times 100$;
    \item the trial time vector, $1 \times 100$;
    \item the earlier recovered embedding, $3 \times 100$;
    \item the delayed recovered embedding, $3 \times 100$.
\end{itemize}

The important point is conceptual rather than file-format specific: the same
trial contains both the generated spike counts and the recovered manifold for
the earlier and delayed populations.

\section{Known Limitations}

The current spike-count observation model is still simplified:

\begin{itemize}
    \item overdispersion, refractoriness, and bursts are phenomenological
    corrections rather than mechanistic biophysical models;
    \item shared population noise and neuron-neuron correlations are still not
    explicitly modeled;
    \item neuron-specific spike variability is limited;
    \item refractory and burst parameters are currently simple distributional
    controls rather than fitted parameters.
\end{itemize}

\section{Selected References}

\begin{itemize}
    \item Dayan, P. and Abbott, L. F. (2001).
    \textit{Theoretical Neuroscience: Computational and Mathematical Modeling
    of Neural Systems}. MIT Press.

    \item Jolliffe, I. T. (2002).
    \textit{Principal Component Analysis}. Springer.

    \item LeCun, Y., Bengio, Y., and Hinton, G. (2015).
    Deep learning.
    \textit{Nature}, 521, 436--444.
    \texttt{https://doi.org/10.1038/nature14539}

    \item Vaswani, A., et al. (2017).
    Attention is all you need.
    \textit{Advances in Neural Information Processing Systems}.
    \texttt{https://arxiv.org/abs/1706.03762}

    \item van den Oord, A., Li, Y., and Vinyals, O. (2018).
    Representation learning with contrastive predictive coding.
    \texttt{https://arxiv.org/abs/1807.03748}

    \item Gallego, J. A., et al. (2023).
    Aligning latent representations of neural activity.
    \textit{Nature Biomedical Engineering}, 7, 337--343.
    \texttt{https://doi.org/10.1038/s41551-022-00962-7}

    \item Pezzulo, G., Donnarumma, F., Ferrari-Toniolo, S., Cisek, P.,
    and Battaglia-Mayer, A. (2022).
    Shared population-level dynamics in monkey premotor cortex during solo
    action, joint action and action observation.
    \textit{Progress in Neurobiology}, 210, 102214.
    \texttt{https://doi.org/10.1016/j.pneurobio.2021.102214}

    \item Schoenemann, P. H. (1966).
    A generalized solution of the orthogonal Procrustes problem.
    \textit{Psychometrika}, 31, 1--10.
    \texttt{https://doi.org/10.1007/BF02289451}

    \item Kriegeskorte, N., Mur, M., and Bandettini, P. (2008).
    Representational similarity analysis: connecting the branches of systems
    neuroscience.
    \textit{Frontiers in Systems Neuroscience}, 2, 4.
    \texttt{https://doi.org/10.3389/neuro.06.004.2008}

    \item Haxby, J. V., et al. (2011).
    A common, high-dimensional model of the representational space in human
    ventral temporal cortex.
    \textit{Neuron}, 72, 404--416.
    \texttt{https://doi.org/10.1016/j.neuron.2011.08.026}

    \item Kornblith, S., Norouzi, M., Lee, H., and Hinton, G. (2019).
    Similarity of neural network representations revisited.
    \textit{Proceedings of ICML}.
    \texttt{https://arxiv.org/abs/1905.00414}

\end{itemize}

\end{document}
